\section{Continue and verify: regime semantics}
\label{sec:b-regimes}

The first study's action set has two members.  Their semantics are fixed
before any data are collected; a change to either is a new action version and
invalidates the qualification of the affected branches.

\begin{table}[t]
\centering
\small
\setlength{\tabcolsep}{4pt}
\begin{tabular}{@{}p{0.14\linewidth}p{0.40\linewidth}p{0.38\linewidth}@{}}
\toprule
Regime & Action & Termination \\
\midrule
$g_0$ & Continue with the frozen agent policy $\rho$; no external call. & Runs to the registered horizon; the root's ordinary outcome. \\
$g_v$ & Invoke the fixed verifier on the declared verifier input, apply the fixed result-to-response rule, then resume $\rho$. & The response rule either leaves the candidate unchanged or replaces it according to the rule; then $\rho$ runs to the horizon. \\
\bottomrule
\end{tabular}
\caption{Regime semantics for the first study.  Verification that only logs a
diagnosis without changing the continuation is a different treatment and is not
$g_v$.}
\label{tab:option-boundaries}
\end{table}

\paragraph{The verifier contract.}
$g_v$ is defined by a frozen tuple: the verifier's input (the declared state
surface and the candidate artifact), its model/tool binding and version, its
output schema, its error handling, the result-to-response rule that maps a
verifier verdict to a change in the continuation, its maximum work (token,
time and call ceilings), its termination condition, and the point at which
control returns to $\rho$.  ``Independent verifier'' is used only in the defined
sense of \emph{separately invoked and unable to access hidden evaluation
labels}; independence of errors is not inferred from a different prompt or model
name, and no verifier is described as independent without those properties.

\paragraph{Feasibility.}
$A_{ig_v}$ is determined before the intervention from the pre-treatment state
and recorded with its reason and evidence.  Three cases are kept distinct:
\begin{itemize}[leftmargin=1.4em]
\item \emph{Unavailable}: the verifier's declared input is absent, so no
  counterfactual outcome is assigned under that action definition.  The
  checkpoint contributes to the feasible-support analysis only through the
  reported unavailable stratum.
\item \emph{Feasible but the verifier times out or errors}: the branch remains
  an outcome of $g_v$, with the execution status and failure reason recorded.
\item \emph{Evaluator or collection failure}: retained as a separate
  data-integrity status and handled by the predeclared missing-data rule; it is
  not interchangeable with task failure.
\end{itemize}
A controller asked to use an unavailable action falls back to $g_0$ explicitly,
and the attempted-invalid-action event is reported.

\paragraph{Recovery is optional and not in the first study.}
A recovery action $g_r$ enters only with a stable precondition, a declared
rollback boundary, a repair action and a continuation contract; a generic
unconstrained ``try again'' is not a comparable regime.  The first study
therefore reports continue/verify only, and \cref{sec:b-bridge} states the
increment a qualified $g_r$ would have to demonstrate.
